%% =====================================================================
%% WS-3 Reference Method Paper — MICCAI 2027 submission
%%
%% Target venue: MICCAI 2027 (Lecture Notes in Computer Science, llncs.cls)
%% Fallback:     IEEE Transactions on Medical Imaging
%% Format:       LNCS, 8 pages main + up to 4 pages supplementary
%% Template:     Springer LNCS. Class file llncs.cls available from
%%               https://www.springer.com/gp/computer-science/lncs
%%               (we use a vanilla \documentclass[runningheads]{llncs};
%%               replace with the real cls at submission).
%%
%% Status: v0.1 working draft. RESULTS tables / figures are \todo{}
%%         placeholders until WS-3 Phase 3 (v1 with ensembles + cross-vendor)
%%         lands at D9+270.
%%
%% Build:  pdflatex manuscript && bibtex manuscript && pdflatex manuscript
%% =====================================================================

\documentclass{article}  % swap for `\documentclass[runningheads]{llncs}` at submission

\usepackage[utf8]{inputenc}
\usepackage[T1]{fontenc}
\usepackage{lmodern}
\usepackage{microtype}
\usepackage[margin=1in]{geometry}
\usepackage{graphicx}
\usepackage{booktabs}
\usepackage{amsmath,amssymb}
\usepackage{xfrac}
\usepackage{algorithm}
\usepackage{algpseudocode}
\usepackage{xcolor}
\usepackage[colorlinks=true,linkcolor=blue,citecolor=blue,urlcolor=blue]{hyperref}
\usepackage{xurl}
\usepackage[round]{natbib}
\usepackage{authblk}

\newcommand{\todo}[1]{\textcolor{red}{\textbf{[TODO: #1]}}}
\newcommand{\R}{\mathbb{R}}
\newcommand{\todoTable}[2]{%
  \begin{table}[ht]\centering%
    \fbox{\parbox{0.85\linewidth}{\centering\todo{TABLE: #2}}}%
    \caption{#2 (placeholder; populate when data lands).}\label{#1}%
  \end{table}}
\newcommand{\todoFigure}[2]{%
  \begin{figure}[ht]\centering%
    \fbox{\parbox{0.85\linewidth}{\centering\todo{FIGURE: #2}}}%
    \caption{#2 (placeholder; populate when data lands).}\label{#1}%
  \end{figure}}

% Stub the LNCS-only \keywords macro so the draft compiles under `article`.
% The real llncs.cls supplies this; remove the next line at submission.
\providecommand{\keywords}[1]{\par\noindent\textbf{Keywords:} #1\par}

\title{\textbf{Unrolled Iterative Reconstruction with Calibrated Uncertainty for Low-Dose CT: \\ A Reference Method for the PWM Benchmark}}

\author[1,2]{\textit{Authors to be confirmed at submission}}
\affil[1]{University of Texas Southwestern Medical Center, Dallas, TX, USA}
\affil[2]{PWM Protocol Foundation}

\date{Working draft v0.1 --- \today}

\begin{document}
\maketitle

%% =====================================================================
\begin{abstract}
We present an open-source reference reconstruction method for low-dose CT designed as a \textbf{credible, improvable, and permissively-licensed} community baseline rather than a single-dataset state-of-the-art point. The method is an unrolled iterative reconstruction with $10$ data-consistency steps interleaved with a weight-shared U-Net denoiser; uncertainty is quantified via a five-model deep ensemble; the forward model is the open \texttt{ct\_radon} implementation maintained as part of the PWM modality engine, so any user replicating the architecture inherits a Radon transform validated against PWM's 172-modality test suite. We evaluate on the multi-vendor PWM-LDCT benchmark across four dose levels (10\%, 25\%, 50\%, 100\%) and report five complementary axes that are rarely co-located in published methods: reconstruction quality (PSNR / SSIM / LPIPS), downstream task performance (lung-nodule detection AUC), cross-vendor generalization (leave-one-vendor-out), per-pixel uncertainty calibration (Spearman $\rho$ of predicted variance vs observed error), and signal-equivalence credentials per the framework of \citet{ws2framework}. The method ships under Apache~2.0 with a content-addressed Docker RunBundle on IPFS and a verified L4 cert on the PWM blockchain registry. Its architecture is \emph{decomposably improvable}: each of the four components (denoiser, forward model, iteration schedule, UQ method) can be swapped independently, inviting follow-on contributions to compete on individual axes rather than re-derive the full stack.
\end{abstract}

\keywords{Low-dose CT \and Unrolled iterative reconstruction \and Uncertainty quantification \and Cross-vendor generalization \and Reproducible benchmark.}

%% =====================================================================
\section{Introduction}
\label{sec:intro}

Despite years of method papers reporting fractional-dB PSNR gains over previous best~\citep{chen2017redcnn,wang2020deeplearningreview,liu2024diffusion}, the low-dose CT field has no single reference baseline. The obstacle is structural: published methods optimize for performance on the dataset used during development (typically AAPM-2016~\citep{aapm2016challenge} on Siemens hardware) and rarely report \textbf{(i)} cross-vendor generalization, \textbf{(ii)} per-pixel uncertainty calibration, \textbf{(iii)} downstream task performance, or \textbf{(iv)} open-source code with permissive licensing. A deployer picking a method for a different vendor, or one needing per-case diagnostic confidence, has no clear choice.

We present an open-source reference method explicitly designed for these four properties, positioned as the seed entry of a permanent benchmark leaderboard~\citep{ws4leaderboard}.

\paragraph{Position vs vendor IR.} Vendor-deployed iterative reconstruction (Siemens \textit{Admire}, GE \textit{Veo}, Canon \textit{AIDR3D}) is FDA-cleared but closed: weights, training data, and inference code are vendor-proprietary; cross-vendor evaluation is therefore not possible. Open methods such as RED-CNN~\citep{chen2017redcnn} are reproducible but lack UQ and cross-vendor scope. We aim to combine the architectural class of vendor IR with the reproducibility of academic methods.

\paragraph{Architecture rationale.} We use unrolled iterative reconstruction~\citep{adler2018learned,mardani2018neural} for three reasons. \textbf{(1) Physical prior:} embedding the (exact) Radon operator in the iteration loop guarantees data consistency that end-to-end networks must learn from data. \textbf{(2) Decomposable improvability:} a follow-up can independently swap the denoiser, forward model, iteration schedule, or UQ method --- each yielding a publishable contribution that improves on this baseline without invalidating it. Pure transformer or pure diffusion architectures do not offer this decomposability~\citep{liu2024diffusion}. \textbf{(3) Regulatory transferability:} iterative methods with explicit data-consistency loops have FDA-cleared precedent in vendor IR; an architecturally analogous open method is more transferable to regulated deployment than a pure neural approach.

\paragraph{Why ``credible and improvable'' beats ``state-of-the-art''.} A reference method that wins by 0.3~dB on a single vendor is too fragile to anchor a community benchmark: the next year's method on a slightly different dataset displaces it, and the cumulative leaderboard has no stable origin. We instead target a competitive, mid-range performer with four desirable properties present, designed for follow-on improvement on independent axes. This is the same logic that made BERT-base a productive baseline relative to a hypothetical singularly-tuned competitor; the open question of which axis to improve next is itself a contribution to the field.

\paragraph{Contributions.}
\begin{enumerate}
    \item An open-source unrolled iterative architecture with weight-shared U-Net denoiser, TV warm start, and five-model deep ensemble UQ, sized to fit in $< 24$~GB GPU memory (Section~\ref{sec:method}).
    \item Cross-vendor generalization evaluation: leave-one-vendor-out on the PWM-LDCT benchmark~\citep{ws1dataset} (Section~\ref{sec:experiments}).
    \item Per-pixel uncertainty calibration: Spearman correlation between predicted variance and observed reconstruction error (Section~\ref{sec:results}).
    \item Five-tuple signal-equivalence credentials per~\citet{ws2framework} at 25\% and 10\% dose (Section~\ref{sec:credentials}).
    \item Release under Apache~2.0 with a reproducible Docker RunBundle on IPFS and a corresponding L4 cert on PWM mainnet.
\end{enumerate}

%% =====================================================================
\section{Method}
\label{sec:method}

\subsection{Forward model}

Let $\mathcal{R}: \R^{H \times W} \to \R^{n_{\textrm{views}} \times n_{\textrm{dets}}}$ denote the Radon forward operator and $\mathcal{R}^\dagger$ its filtered back-projection adjoint. Both are implemented in \texttt{pwm\_core.contrib.modalities.ct\_radon}~\citep{pwmcore}, the validated forward-model implementation used across PWM's 172-modality engine. We do not re-implement; using the upstream implementation guarantees that bug-fixes propagate uniformly and that any user of our reference method inherits a Radon transform validated against the full PWM modality test suite.

\subsection{Unrolled iteration}

Given a low-dose sinogram $y \in \R^{n_{\textrm{views}} \times n_{\textrm{dets}}}$, we estimate the reconstruction $\hat{x} \in \R^{H \times W}$ via $K$ iterations of a data-consistency step interleaved with a learned denoiser $f_\theta$:

\begin{algorithm}[h]
    \caption{Unrolled iterative reconstruction}
    \label{alg:unrolled}
    \begin{algorithmic}[1]
        \State $x^{(0)} \gets \textrm{TV-FBP}(y)$ \Comment{TV-regularized warm start}
        \For{$k = 0, 1, \ldots, K{-}1$}
            \State $x^{(k+\sfrac{1}{2})} \gets x^{(k)} - \tau \cdot \mathcal{R}^\dagger \!\left( \mathcal{R}(x^{(k)}) - y \right)$ \Comment{Data consistency}
            \State $x^{(k+1)} \gets f_\theta\!\left( x^{(k+\sfrac{1}{2})} \right)$ \Comment{Learned denoising}
        \EndFor
        \State \Return $\hat{x} \gets x^{(K)}$
    \end{algorithmic}
\end{algorithm}

We use $K = 10$ iterations; ablation results for $K \in \{5, 10, 15, 20\}$ are in Section~\ref{sec:results}~Table~\ref{tab:K_ablation}. The denoiser $f_\theta$ is a U-Net with $\approx 4$~M parameters whose weights are shared across iterations; this both reduces memory and acts as an implicit regularizer (the same $f_\theta$ must denoise at all iteration depths, preventing iteration-specific overfitting).

\todoTable{tab:K_ablation}{Iteration-count ablation: PSNR / SSIM on PWM-LDCT test split for $K \in \{5, 10, 15, 20\}$}

The step size $\tau$ is a learned scalar, initialized at $0.5 / \|\mathcal{R}\|^2$ (the conventional Landweber-stable value). End-to-end optimization with mean absolute error loss on Hounsfield-unit values:

\begin{equation}
    \mathcal{L}(\theta, \tau) = \frac{1}{N}\sum_{i=1}^N \left\| \hat{x}_i(\theta, \tau; y_i) - x_i^{\textrm{ref}} \right\|_1
\end{equation}

where $x_i^{\textrm{ref}}$ is the patient's full-dose reconstruction (the reference image).

\subsection{Uncertainty quantification}

We use a five-model deep ensemble~\citep{lakshminarayanan2017deep}: five independently-initialized and independently-trained instances of the same architecture, differing only in random seed for weight initialization and data ordering. At inference time, given a low-dose sinogram $y$, each ensemble member produces a reconstruction $\hat{x}_m(y)$ for $m = 1, \ldots, 5$. The point estimate is the mean and the per-pixel uncertainty is the standard deviation:

\begin{equation}
    \hat{x}(y) = \frac{1}{5}\sum_{m=1}^{5} \hat{x}_m(y), \qquad \sigma(y)_{ij} = \sqrt{\frac{1}{5}\sum_{m=1}^{5}\left(\hat{x}_m(y)_{ij} - \hat{x}(y)_{ij}\right)^2}.
\end{equation}

We choose deep ensembles over Monte-Carlo dropout because (i) ensembles consistently outperform dropout on calibration in regression benchmarks~\citep{ovadia2019uncertainty}, (ii) inference cost is bounded ($5\times$ a single model rather than $T$-stochastic), and (iii) the inference cost can be amortized across many predictions on the same GPU.

\paragraph{What the uncertainty is \emph{for}.} A clinically actionable UQ map distinguishes patient cases for which the reconstruction is reliable (low $\sigma(y)$ everywhere in the diagnostic region) from cases for which the reconstruction is uncertain enough to warrant either a higher-dose retake, a second-reader review, or a deferred diagnosis. A well-calibrated reference method therefore reports not only per-pixel variance but its \emph{agreement with per-pixel error} on a held-out test set (Section~\ref{sec:results}). A high Spearman correlation gives radiologists evidence that the variance map is informative; a low correlation flags that the ensemble disagreement does not track the actual reconstruction quality and the UQ should not be used for triage. The PWM-LDCT benchmark's paired full-dose ground truth makes this evaluation possible at scale.

\subsection{Training data and augmentation}

Training data: full PWM-LDCT~\citep{ws1dataset} training split with paired (full-dose, low-dose) reconstructions at $r \in \{0.10, 0.25, 0.50\}$. We mix all three dose levels in each minibatch with weights $(0.4, 0.4, 0.2)$, biased toward the most clinically relevant levels.

Augmentation: random flip ($p=0.5$), random rotation in $[-5^\circ, 5^\circ]$, random window-level shift in $[-50, 50]$~HU. No augmentation that violates physical consistency (no random noise injection on top of the already-noisy low-dose, no patch shuffling).

%% =====================================================================
\section{Experiments}
\label{sec:experiments}

\paragraph{Datasets.} Training: PWM-LDCT~\citep{ws1dataset} training split (\todo{N} patients, \todo{N} scans). Evaluation: PWM-LDCT test split (\todo{N} patients), AAPM-2016~\citep{aapm2016challenge} held-out scans (\todo{N} patients), and LIDC-IDRI~\citep{lidc_idri} sub-sample (\todo{N} patients) for downstream task validation.

\paragraph{Baselines.} Three baseline reconstruction methods, reproduced within $\leq 0.5$~dB PSNR of their published numbers:
\begin{itemize}
    \item RED-CNN~\citep{chen2017redcnn}: foundational encoder-decoder CNN, 2017.
    \item A transformer-based reconstruction (\todo{select: CTformer / TransCT / Restormer-LDCT}).
    \item A diffusion / score-based reconstruction (\todo{select: 2024 score-LDCT}).
\end{itemize}

\paragraph{Cross-vendor splits.} We construct three leave-one-vendor-out splits: train-on-\{GE, Canon\}/test-on-Siemens; train-on-\{Siemens, Canon\}/test-on-GE; train-on-\{Siemens, GE\}/test-on-Canon. Each cross-vendor experiment uses the full training set with the held-out vendor's patients removed (patient-level, not slice-level).

\paragraph{Metrics.} (i) PSNR, SSIM, LPIPS~\citep{zhang2018lpips} on the test reconstructions; (ii) lung-nodule detection AUC at fixed FPR = 0.1, using a pre-trained nodule detector frozen across reconstruction methods; (iii) per-pixel uncertainty calibration --- Spearman correlation of $\sigma(y)_{ij}$ against $|\hat{x}(y)_{ij} - x_i^{\textrm{ref}}|_{ij}$ on held-out test; (iv) inference latency (seconds per slice on a single NVIDIA A100), peak GPU memory; (v) 5-tuple signal-equivalence credentials per the framework of \citet{ws2framework} at $r \in \{0.10, 0.25\}$ for lung-nodule detection, $\varepsilon = 0.02$, $\alpha = 0.05$, subpopulation $\Pi = $ adult-chest test split.

%% =====================================================================
\section{Results}
\label{sec:results}

\paragraph{Reconstruction quality.} Table~\ref{tab:reconstruction_quality} \todo{populate} reports PSNR, SSIM, and LPIPS for each method at each dose level on the PWM-LDCT test split. The reference method achieves PSNR within $\todo{X}$~dB of the diffusion-based baseline at $r = 0.25$, while consuming $\todo{Y}\times$ less inference time and using $\todo{Z}\times$ less GPU memory.

\begin{table}[h]
    \centering
    \small
    \setlength{\tabcolsep}{4pt}
    \begin{tabular}{lcccccc}
        \toprule
        \textbf{Method} & \multicolumn{2}{c}{$r = 0.10$} & \multicolumn{2}{c}{$r = 0.25$} & \multicolumn{2}{c}{$r = 0.50$} \\
        & PSNR & SSIM & PSNR & SSIM & PSNR & SSIM \\
        \midrule
        FBP                & \todo{} & \todo{} & \todo{} & \todo{} & \todo{} & \todo{} \\
        RED-CNN            & \todo{} & \todo{} & \todo{} & \todo{} & \todo{} & \todo{} \\
        Transformer        & \todo{} & \todo{} & \todo{} & \todo{} & \todo{} & \todo{} \\
        Diffusion          & \todo{} & \todo{} & \todo{} & \todo{} & \todo{} & \todo{} \\
        \textbf{PWM ref}   & \todo{} & \todo{} & \todo{} & \todo{} & \todo{} & \todo{} \\
        \bottomrule
    \end{tabular}
    \caption{Reconstruction quality on the PWM-LDCT test split. Values to fill at D9+270 after WS-3 Phase 3 (v1 deep ensemble) lands.}
    \label{tab:reconstruction_quality}
\end{table}

\paragraph{Cross-vendor generalization.} Table~\ref{tab:cross_vendor} reports the PSNR drop $\Delta$ when each method is trained without one vendor and evaluated on that held-out vendor's test scans. The reference method's $\Delta$ should be $\leq 2$~dB; baselines are not designed for cross-vendor generalization and we expect $\Delta \approx 3$--$5$~dB.

\todoTable{tab:cross_vendor}{Cross-vendor PSNR drop ($\Delta$, dB) under leave-one-vendor-out training; lower is better}

\paragraph{Uncertainty calibration.} The deep ensemble's per-pixel standard deviation correlates with per-pixel error at Spearman $\rho = $ \todo{populate}. A reliability diagram is provided as Figure~\ref{fig:reliability}.

\todoFigure{fig:reliability}{Reliability diagram: predicted per-pixel standard deviation (ensemble) vs observed per-pixel reconstruction error on held-out test}

\paragraph{Downstream task performance.} Lung-nodule detection AUC against a frozen detector is reported in Table~\ref{tab:task_auc}. The downstream-task results may diverge from pixel-wise results: a method with higher PSNR may produce lower task AUC if it over-smooths small lesions~\citep{kim2019smoothing}. Showing both is the point.

\todoTable{tab:task_auc}{Lung-nodule detection AUC at fixed FPR = 0.1, per method, per dose level, using a frozen pre-trained detector}

\subsection{Signal-equivalence credentials}
\label{sec:credentials}

Per the framework of \citet{ws2framework}, we report 5-tuple credentials for the reference method:

\begin{table}[h]
    \centering
    \footnotesize
    \begin{tabular}{lcccccc}
        \toprule
        $r$ & Task & $\varepsilon$ & $\alpha$ & Subpopulation & Verdict & CI \\
        \midrule
        0.25 & Lung nodule AUC & 0.02 & 0.05 & adult-chest test split  & \todo{} & \todo{} \\
        0.10 & Lung nodule AUC & 0.05 & 0.05 & adult-chest test split  & \todo{} & \todo{} \\
        0.25 & Liver lesion Dice & 0.03 & 0.05 & adult-abdomen test split & \todo{} & \todo{} \\
        \bottomrule
    \end{tabular}
    \caption{Signal-equivalence credentials computed per \citet{ws2framework} against full-dose FBP as the reference. Verdicts: \texttt{PASS} / \texttt{FAIL} / \texttt{INDETERMINATE}.}
    \label{tab:credentials}
\end{table}

These credentials are the canonical artifact that future methods will report against. Reporting only PSNR/SSIM/LPIPS is insufficient; reporting only verdicts loses the CI; reporting both is the standard we propose.

%% =====================================================================
\section{Conclusion}

We have presented an open-source reference reconstruction method for low-dose CT, designed not to maximize PSNR but to be \emph{credible, improvable, and permissively licensed}: cross-vendor evaluated, deep-ensemble UQ-calibrated, signal-equivalence-credentialed per \citep{ws2framework}, and Apache-2.0 licensed. The method is the seed entry of the permanent PWM Low-Dose CT Challenge leaderboard~\citep{ws4leaderboard}; an L4 cert anchoring the RunBundle CID is registered on the PWM blockchain mainnet. We invite the community to compete by improving any of the four independent axes (denoiser, forward model, iteration schedule, UQ method) without re-deriving the full stack. Code, weights, and RunBundle are released under Apache 2.0 at \url{https://github.com/integritynoble/low_dose_CT/tree/main/WS-3_reference_method}.

%% =====================================================================
\section*{Acknowledgments}
\todo{Fill at submission.}

%% =====================================================================
\bibliographystyle{plainnat}
\bibliography{refs}

\end{document}
